\documentclass[11pt]{article}
\usepackage[margin=1in]{geometry}
\usepackage{amsmath,graphicx,booktabs,amssymb,xcolor,hyperref,array}
\usepackage[skip=4pt]{caption}
\hypersetup{colorlinks=true,linkcolor=blue!55!black,urlcolor=blue!55!black}
\newcommand{\fig}[1]{\includegraphics[width=\linewidth]{figs/#1}}
\newcommand{\good}[1]{\textcolor{green!45!black}{#1}}
\newcommand{\bad}[1]{\textcolor{red!70!black}{#1}}
\sloppy
\title{\textbf{When does per-turn credit assignment help?}\\[2pt]
\large A cross-environment study of credit assignment for multi-turn agents\\(AlfWorld $\cdot$ SearchQA)}
\author{Agentic-RL-CA project}
\date{2026-07-26}

\begin{document}
\maketitle

\begin{abstract}
\noindent We compare four credit-assignment schemes for multi-turn RL agents --- outcome-only
group-relative (\textbf{GRPO}), anchor-state group-relative (\textbf{GiGPO}), and two critic-based
per-turn PPO variants (\textbf{turn-PPO}, \textbf{token-PPO}) --- across two environments with very
different horizons: \textbf{AlfWorld} (embodied household tasks, 30--50 turns, Qwen3-1.7B) and
\textbf{SearchQA} (agentic retrieval QA, a hard 4-turn cap, Qwen3-4B). The central finding is that
\emph{the ranking of credit-assignment methods reverses with horizon length}. On long-horizon
AlfWorld the per-turn methods win decisively (turn-PPO 0.963, GiGPO 0.955 vs GRPO 0.881); on
short-horizon SearchQA the flat outcome-only methods win (GRPO 0.441 $\approx$ GiGPO 0.446 vs
turn-PPO 0.389, token-PPO 0.381). A controlled contrast isolates the mechanism: turn-level and
token-level PPO are the same algorithm differing \emph{only} in whether GAE runs over turns or
tokens, and that single change is worth $+0.202$ on AlfWorld but only $+0.008$ on SearchQA --- a
25$\times$ larger effect where the horizon is long. A horizon-stratified analysis explains why:
per-turn credit's value scales with \emph{genuine sequential depth}. On AlfWorld it keeps trajectories short and
on-track (GRPO wastes turns wandering); on SearchQA the 4-turn cap is the binding ceiling, and the
only positive horizon signal is the \emph{critic} methods narrowing GRPO's gap on the deepest
truly-sequential chains (MuSiQue, turn-PPO $+0.031$ [$+0.003,+0.057$]), a benefit swamped by
over-search on shallow comparison questions. A thinking-mode ablation (SearchQA only) confirms the ordering is not a
truncation artifact and exposes a second effect: without thinking, both arms score \emph{higher}
using $\approx$70\% fewer tokens, and GRPO re-spends the freed budget on retrieval --- consuming
3.99 of 4 turns and hitting the cap on 99.5\% of questions, versus 3.07 under thinking --- while
turn-PPO does not (2.70 $\to$ 2.61 turns). We release all figures, the rerunnable analysis, and the recovered required-hops
annotations.
\end{abstract}

\section{Setup}
\textbf{Methods.} \emph{GRPO} --- token-level, outcome-only group-relative advantage ($\gamma{=}1$),
no critic. \emph{GiGPO} --- group-relative with an additional anchor-state step advantage
($\gamma{=}0.95$), no critic. \emph{turn-PPO} --- learned critic + cross-turn GAE
(\texttt{gae\_turn}, $\gamma{=}\lambda{=}0.95$). \emph{token-PPO} --- learned critic + token-level
GAE (\texttt{gae}); the granularity control for turn-PPO. \emph{floor} --- the pre-RL policy.

\textbf{Environments \& scale (a caveat carried throughout).} AlfWorld uses Qwen3-\textbf{1.7B}
(reward\_models campaign, single seed, RL from a replay-BC policy); SearchQA uses Qwen3-\textbf{4B}
(this project, seed 0). The two are different environments and not scale-matched --- cross-environment
claims are about the \emph{pattern of method rankings}, never absolute numbers. All four arms now run in both
environments (the AlfWorld token-PPO leg was added 2026-07-24 to complete the granularity contrast).

\textbf{Metric \& protocol.} AlfWorld: task success rate (seen 128-game greedy val for training
curves; unseen 134-game greedy for eval). SearchQA: per-question exact match (EM) with alias
normalisation over the full 7-task, 51{,}713-question test set (NQ, TriviaQA, PopQA, HotpotQA,
2WikiMultihopQA, MuSiQue, Bamboogle), greedy, top-3 wiki-18 retrieval, 4-turn cap. Identical harness
for every checkpoint.

\textbf{Required hops ($H_{\text{gold}}$).} The SearchQA packing had stripped the datasets to
question+answer; we re-attached the annotated reasoning-chain length by an exact positional join to
the original FlashRAG splits (verified $51{,}713/51{,}713$ question match). $H_{\text{gold}}{=}1$ for
NQ/TriviaQA/PopQA, $2$ for HotpotQA/Bamboogle, $\text{len(decomposition)}\in\{2,3,4\}$ for MuSiQue,
and $2$ or $4$ for 2Wiki (bridge\_comparison). Strata sizes: $H1{=}29{,}190$, $H2{=}18{,}607$,
$H3{=}760$, $H4{=}3{,}156$.

\section{Headline: the method ranking reverses across environments}

\begin{table}[h]\centering\footnotesize\setlength{\tabcolsep}{4.2pt}
\caption{Evaluation accuracy. AlfWorld = unseen success rate; SearchQA = full-set macro-EM (7-task
mean). Best per column \good{green}, worst trained arm \bad{red}. Two things to read: (i) the
\emph{reversal} --- turn-PPO beats GRPO on AlfWorld (0.963 vs 0.881) while GRPO beats turn-PPO on
SearchQA (0.441 vs 0.389); (ii) \emph{token}-PPO is worst in both environments, and catastrophically
so on long-horizon AlfWorld (0.761), where it differs from turn-PPO \emph{only} in GAE granularity.
SearchQA columns are the \textbf{full 51{,}713-question} test set (Fig 1 plots the 2{,}048-question
training proxy). \textbf{GRPO and GiGPO are tied}: GiGPO leads by $0.005$ here and trails by $0.007$
on the proxy --- both below the proxy's own $\sim$0.01 error, so no ordering between them is claimed.}
\begin{tabular}{lcccccccc c}
\toprule
& \multicolumn{1}{c}{\textbf{AlfWorld}} & \multicolumn{8}{c}{\textbf{SearchQA (per-task EM \& macro)}}\\
\cmidrule(lr){2-2}\cmidrule(lr){3-10}
Method & unseen & NQ & TriviaQA & PopQA & HotpotQA & 2Wiki & MuSiQue & Bamb. & \textbf{macro}\\
\midrule
pre-RL floor & 0.440 & 0.292 & 0.517 & 0.346 & 0.271 & 0.331 & 0.074 & 0.360 & 0.313\\
token-PPO    & \bad{0.761} & 0.454 & 0.595 & 0.447 & 0.365 & 0.323 & 0.137 & 0.344 & \bad{0.381}\\
turn-PPO     & \good{0.963} & 0.457 & 0.620 & 0.450 & 0.380 & 0.338 & 0.128 & 0.352 & 0.389\\
GRPO         & 0.881        & 0.483 & 0.643 & 0.478 & 0.435 & 0.449 & 0.184 & 0.416 & 0.441\\
GiGPO        & 0.955 & 0.490 & 0.647 & 0.469 & 0.447 & 0.399 & 0.187 & 0.480 & \good{0.446}\\
\bottomrule
\end{tabular}
\end{table}

\begin{figure}[h]\centering\fig{fig1_training_curves.pdf}
\caption{\textbf{Fig 1.} Training accuracy. \emph{Left, AlfWorld:} seen 128-game success --- the
per-turn methods (GiGPO 0.992, turn-PPO 0.961) dominate flat GRPO (0.930). \emph{Right, SearchQA:}
\textbf{val-2048 macro-EM} --- the training-time proxy on a 2{,}048-question subset, \emph{not} the
full-set metric of Table 1. The order inverts relative to AlfWorld: group-relative (GRPO 0.443
$\approx$ GiGPO 0.436) beats the critic-PPO arms (turn 0.389, token 0.364). \textbf{NB} GRPO and
GiGPO swap places between this proxy ($-0.007$) and the full-set eval ($+0.005$); the proxy tracks
GRPO to $0.002$ but under-reads GiGPO by $0.010$, so the two are best read as \emph{tied} (see
Table 1 note). (AlfWorld PPO trained 150 steps vs 300 for GRPO/GiGPO.)}
\end{figure}

GRPO is the \emph{weakest} of the three long-horizon-viable arms on AlfWorld (0.881 vs 0.955/0.963)
yet \emph{among the best} on SearchQA (0.441, statistically tied with GiGPO's 0.446); turn-PPO
makes exactly the opposite trip (best on AlfWorld, 0.052 behind GRPO on SearchQA). Any single-environment conclusion about credit assignment would
therefore be misleading. token-PPO, meanwhile, is worst in both --- but for a reason specific to
granularity, which \S3 isolates. The rest of the report explains the reversal.

\section{Turn behaviour: RL is largely learning \emph{how much to act}}

\begin{figure}[h]\centering
\begin{minipage}{0.46\linewidth}\fig{fig2_alfworld_turns.pdf}\end{minipage}\hfill
\begin{minipage}{0.52\linewidth}\fig{fig3_searchqa_turns_box.pdf}\end{minipage}
\caption{\textbf{Fig 2 (left):} AlfWorld trajectory length --- token-PPO uses \textbf{24.2} turns on
average and GRPO 14.9, vs GiGPO 9.1 / turn-PPO 10.9. \textbf{Fig 3 (right):} SearchQA turns per dataset (cap 4);
turn usage rises with task difficulty and the trained arms sit well above the under-searching floor.}
\end{figure}

\textbf{On SearchQA, RL mostly teaches the model to use its turn budget.} The pre-RL floor
\emph{under-searches} --- median 2 turns / 1 search, macro-EM 0.313. RL roughly doubles search usage
(GRPO/GiGPO: 3 turns / 2 searches) for a $+0.13$ macro lift. But more turns is not monotonically
good: GiGPO \emph{over}-searches into the cap (cap-fail 0.29/0.50/0.88/0.66 across $H1$--$H4$ vs
GRPO's 0.12/0.29/0.72/0.47), which is why it fails to beat flat GRPO despite identical macro-EM.
Turn-calibration, not credit fidelity, is the lever on a 4-turn budget.

\textbf{On AlfWorld, the whole ranking is ``how often does it wander?''} Every method solves
\emph{exactly} 100\% of episodes that finish within $\le 20$ turns (Fig 5, left; all four arms, all
four sub-20 bins). The accuracy differences are entirely about \emph{what fraction of episodes drag
past 20 turns}:

\begin{center}\small
\begin{tabular}{lcccc}
\toprule
& GiGPO & turn-PPO & GRPO & token-PPO\\
\midrule
episodes $>$20 turns & 4.5\% & 11\% & 21\% & \bad{46\%}\\
\quad of which succeed & 0.00\tiny(n=6) & 0.67 & 0.43 & 0.48\\
mean turns & 9.1 & 10.9 & 14.9 & \bad{24.2}\\
unseen success & 0.955 & \good{0.963} & 0.881 & \bad{0.761}\\
\bottomrule
\end{tabular}
\end{center}

A concrete GRPO failure mode, \emph{``examine the book with the desklamp''} solved in 44 turns:
\begin{quote}\ttfamily\small
t0 go to shelf 1 $\;\to\;$ t1 go to shelf 2 $\;\to\;$ t3 go to bed 1 $\;\to\;$ t4 take book 1
$\;\to\;$ t5 go to desk 1 $\;\to\;$ \bad{t6..t13 look / look / look / look / look / look / look / look}
\end{quote}
GRPO loops the no-op \texttt{look} eight times; GiGPO solves the same task type in \textbf{3} turns.
Per-turn credit penalises the wasted steps directly; the outcome-only signal does not.

\subsection*{The cleanest test: turn-level vs token-level GAE}
turn-PPO and token-PPO are the \emph{same algorithm} --- same learned critic, same BC initialisation,
same $\gamma{=}\lambda{=}0.95$, same 150-step budget, same eval checkpoint. They differ in exactly one
thing: whether the GAE recursion runs over \textbf{turns} or over \textbf{tokens}. That single change
is worth:
\begin{center}\small
\begin{tabular}{lccc}
\toprule
& turn-level & token-level & $\Delta$\\
\midrule
AlfWorld (long horizon, $\le$50 turns) & 0.963 & 0.761 & \good{$+0.202$}\\
SearchQA (short horizon, 4-turn cap)   & 0.389 & 0.381 & $+0.008$\\
\bottomrule
\end{tabular}
\end{center}
\noindent A \textbf{25$\times$ larger effect on the long-horizon environment}, from an identical
algorithm differing only in credit granularity. The mechanism is visible in the discount: at
$\gamma{=}0.95$ per \emph{token}, an AlfWorld transcript of $\sim$25k generated tokens discounts early
credit by $0.95^{25000}\!\approx\!0$, so the value signal cannot reach the decisions that matter, and
the policy wanders (46\% of episodes exceed 20 turns; \texttt{cool} collapses to 0.29 success at 43.6
turns). Under a 4-turn SearchQA cap the same discount spans only a few hundred tokens, so
granularity is nearly irrelevant --- and the two arms tie. This is the sharpest single piece of
evidence for the paper's thesis, because it holds everything except granularity fixed.

\section{Thinking vs non-thinking (SearchQA, Qwen3-4B)}
The SearchQA results above use Qwen3's thinking mode at a 2048-token response budget, under which
the critic arm clips 8.4\% of its turns --- and a clipped turn emits no action, wasting one of only
four. That makes truncation a live confound for the central SearchQA claim (flat GRPO $>$ critic
turn-PPO). We therefore re-ran both arms with \texttt{enable\_thinking=False} at the \emph{same}
2048 budget, 500 steps, batch, seed and $\gamma{=}\lambda{=}1$, so \emph{only} the thinking flag
differs. \textbf{This ablation exists for SearchQA only}: AlfWorld was non-thinking throughout (its
transcript prompt could not afford thinking blocks under a 2048-token prompt cap over $\le$50
turns), so there is no thinking counterpart to compare against there.

\begin{table}[h]\centering\small\setlength{\tabcolsep}{5pt}
\caption{Thinking-mode ablation --- full-set metrics (51{,}713 q, step-500 checkpoints; same eval
harness as Table 1). Non-thinking wins on accuracy for \emph{both} methods while generating
$\approx$70\% fewer tokens. The behavioural split is the interesting part: GRPO converts the freed
budget into \emph{search}, turn-PPO does not.}
\begin{tabular}{llccccc}
\toprule
Method & mode & macro-EM & turns & searches & tokens/turn & hits 4-turn cap\\
\midrule
GRPO      & thinking      & 0.441 & 3.07 & 2.07 & 471 & 28\%\\
GRPO      & non-thinking  & \good{\textbf{0.463}} & \good{3.99} & \good{2.99} & \good{111} & 99.5\%\\
\addlinespace
turn-PPO  & thinking      & 0.389 & 2.70 & 1.70 & 696 & 25\%\\
turn-PPO  & non-thinking  & \good{0.417} & 2.61 & 1.61 & \good{213} & 17\%\\
\midrule
\multicolumn{2}{l}{GRPO $-$ turn-PPO} & $0.052\to0.046$ & & & &\\
\bottomrule
\end{tabular}
\end{table}

\textbf{(i) The headline finding is robust.} GRPO still leads turn-PPO once the truncation channel
is largely removed (turn-PPO's clip falls $0.084\to0.057$), and the gap barely moves
($0.052\to0.046$). The ordering is not a clipping artifact.

\textbf{(ii) Thinking does not pay for itself on this task.} Both methods score \emph{higher}
without it ($+0.022$ GRPO, $+0.028$ turn-PPO) at $\approx$70\% fewer generated tokens. With only
four turns and a retrieval tool available, the budget is better spent acting than deliberating.

\textbf{(iii) Thinking trades search calls for deliberation tokens --- and search wins.} This is
the sharpest read. Non-thinking GRPO learns to consume its \emph{entire} turn budget: 3.99 of 4
turns and $\approx$3 searches, hitting the cap on \textbf{99.5\%} of questions, versus 3.07 turns
under thinking. The tokens freed by dropping the \texttt{<think>} block are re-spent on retrieval,
and that is where the $+0.022$ comes from. \textbf{turn-PPO fails to make the same trade}: its turn
count is essentially unchanged ($2.70\to2.61$), so it banks the token savings without converting
them into extra evidence --- a plausible reason GRPO retains its lead. Note this is the same
turn-under-use that \S4 diagnoses on the hop axis, and it is \emph{not} a truncation effect: at
non-thinking lengths ($\approx$200 tokens/turn) neither arm is near the cap.

\emph{Caveats:} single seed per cell, so the $+0.022$/$+0.028$ margins are suggestive rather than
established; GRPO's advantage in particular sits close to its own run-to-run oscillation band.

\section{Horizon-stratified analysis}

\begin{figure}[h]\centering\fig{fig5_acc_vs_turns.pdf}
\caption{\textbf{Fig 5.} Accuracy vs turns \emph{used}. \emph{Left, AlfWorld:} all methods succeed
until $>20$ turns, where the wandering tail (largest for GRPO) fails. \emph{Right, SearchQA:}
accuracy falls with turns used --- long trajectories are the hard, cap-hitting ones --- and the
critic-PPO arms sit below GRPO/GiGPO at every turn count.}
\end{figure}

\begin{figure}[h]\centering\fig{fig4_searchqa_turns_vs_gold.pdf}
\caption{\textbf{Fig 4.} SearchQA turns used vs required hops, split accurate/failed/all (band
$\pm1$SD; dotted $y{=}x$). All methods escalate turns with hops on MuSiQue and pin at the cap by
$H{=}3$; failed trajectories sit at the cap (budget exhausted) while successes finish earlier.}
\end{figure}

\begin{minipage}{0.5\linewidth}
\textbf{EM vs required hops (within task).} MuSiQue collapses monotonically for every method.
On 2Wiki the floor and GRPO \emph{rise} at the 4-hop (bridge\_comparison) stratum --- the
comparison/base-rate shortcut --- while the over-searching arms fall, and the critic arms end
\emph{below the untrained floor} there (0.224/0.138 vs 0.456): learned search displaces the
floor's successful parametric early answers. The MuSiQue--2Wiki gap is the crux of \S\ref{sec:why}.
\end{minipage}\hfill
\begin{minipage}{0.46\linewidth}\centering\footnotesize\setlength{\tabcolsep}{4pt}
\captionof{table}{Within-task EM by $H_{\text{gold}}$.}
\begin{tabular}{lccc c cc}
\toprule
& \multicolumn{3}{c}{MuSiQue} && \multicolumn{2}{c}{2Wiki}\\
\cmidrule(lr){2-4}\cmidrule(lr){6-7}
& H2 & H3 & H4 && H2 & H4\\
\midrule
floor    & .110 & .042 & .025 && .297 & .456\\
GRPO     & .268 & .114 & .054 && .433 & .504\\
GiGPO    & .276 & .114 & .047 && .406 & .372\\
turn-PPO & .197 & .064 & .035 && .370 & .224\\
token-PPO & .217 & .063 & .025 && .374 & .138\\
\bottomrule
\end{tabular}
\end{minipage}

\medskip
\textbf{Horizon-gain vs GRPO (A2).} We measure, for each arm $m$, the top-minus-bottom contrast of
$\Delta_m(h){=}\text{EM}_m(h){-}\text{EM}_{\text{GRPO}}(h)$ with a 1000-resample item bootstrap; a
horizon effect is claimed only when the CI excludes 0.

\begin{table}[h]\centering\small
\caption{Horizon-gain contrasts vs GRPO ($\Delta(\text{high}){-}\Delta(\text{low})$, 95\% CI).
$^\star$ = CI excludes 0. The \emph{only} positive signal is the critic methods on MuSiQue's deep
chains; every arm loses on 2Wiki by over-searching shallow comparison questions.}
\begin{tabular}{lccc}
\toprule
Method & within-MuSiQue & within-2Wiki & pooled\\
\midrule
turn-PPO & \good{$+0.031$ [$+0.003,+0.057$]$^\star$} & \bad{$-0.217^\star$} & $-0.183^\star$\\
token-PPO & $+0.007$ [$-0.018,+0.032$] & \bad{$-0.307^\star$} & $-0.233^\star$\\
GiGPO & $-0.011$ [$-0.035,+0.013$] & \bad{$-0.106^\star$} & $-0.092^\star$\\
\bottomrule
\end{tabular}
\end{table}

Reading: turn-PPO is worse than GRPO at \emph{every} MuSiQue hop, but its gap \emph{narrows} as hops
deepen ($\Delta$: $-0.071\to-0.019$) --- a small, significant positive contrast. The critic's per-turn
value estimates do buy something on genuinely-sequential chains; it is simply overwhelmed by
over-search losses elsewhere and by lower overall performance. GiGPO's group-relative step advantage
shows no such MuSiQue signal.

\section{Why hop-depth crushes MuSiQue but not 2Wiki}\label{sec:why}
MuSiQue's hop count is genuine \emph{sequential} depth; 2Wiki's ``4-hop'' (bridge\_comparison) is a
\emph{shallow parallel comparison} the label over-counts. A MuSiQue 4-hop is a chain in which each
hop's answer feeds the next (the annotation uses \texttt{\#1/\#2/\#3} placeholders):
\begin{quote}\small
\emph{``president of the newly-independent country that established the Timor-Leste Commission \dots
with the country containing the airport that includes Lion Air?''}\\
airport of Lion Air $\to$ \texttt{Juanda} $\to$ its country $\to$ \texttt{Indonesia} $\to$ +Timor-Leste
Commission $\to$ \texttt{East Timor} $\to$ its president $\to$ \texttt{Francisco Guterres}.
\end{quote}
You cannot begin hop 4 until hops 1--3 resolve. With a 4-turn budget ($\approx$3 searches), the model
runs out mid-chain --- here GRPO is still decomposing at the cap and scores 0:
\begin{quote}\ttfamily\scriptsize
turns=4, em=0: ``<think> \dots the length of the city where the Yongle Emperor greeted a person the
Ming court thought were representatives sent by \dots Yaxing Coach's headquarters. First, I need to
break do[wn]\dots'' \bad{[budget exhausted]}
\end{quote}
A 2Wiki bridge\_comparison, by contrast, is two \emph{independent} 2-hop facts joined by a comparison
(\emph{``which film's director was born later, A or B?''}): the branches don't depend on each other
(depth 2), the answer is one of two named entities ($\sim$50\% base rate), and parametric knowledge
often supplies one side. Hence (i) MuSiQue is \emph{cap-limited} --- cap-fail 0.72 at $H3$, and no
credit scheme overcomes ``4 turns can't do 4 hops''; (ii) 2Wiki-4hop is \emph{easier} than
2Wiki-2hop for every method, floor included, so the difficulty structure is a property of the tasks,
not of the credit signal. This is the composition confound that makes pooled cross-stratum curves
misleading and the within-task curves the honest ones.

\section{Per-method verdict}
\begin{itemize}\itemsep2pt
\item\textbf{GRPO} --- best flat baseline. Tied for best on SearchQA with GiGPO (cheap, no critic,
well-calibrated turn use; 0.441 vs 0.446 full-set, ordering reverses on the proxy)
but wastes the loose AlfWorld budget wandering (14.9 turns; \texttt{look}-loops), making it the
\emph{worst} long-horizon arm.
\item\textbf{GiGPO} --- best long-horizon efficiency (AlfWorld 0.955 at 9.1 turns). On SearchQA it
\emph{ties} GRPO on EM but over-searches into the cap and gains nothing on the hop axis --- its edge
is variance/stability, not credit fidelity.
\item\textbf{turn-PPO} --- best on AlfWorld (0.963). On SearchQA it underperforms overall yet is the
\emph{only} arm with a significant positive horizon-gain, specifically on MuSiQue's deep chains ---
the clearest evidence that critic-based per-turn credit helps in proportion to sequential depth.
\item\textbf{token-PPO} --- worst arm in \emph{both} environments, and the informative failure.
On SearchQA (0.381) it merely trails, with the same MuSiQue-leaning / 2Wiki-losing shape as turn-PPO.
On AlfWorld it collapses to 0.761 at 24.2 turns/episode --- 46\% of episodes exceed 20 turns --- because
per-token discounting annihilates credit over a 25k-token transcript. Identical to turn-PPO in every
respect but granularity, it is the control that turns the paper's thesis from correlation into a
controlled contrast.
\end{itemize}

\section{Discussion \& limitations}
\textbf{Thesis.} The value of fine-grained per-turn credit assignment scales with \emph{genuine
sequential horizon}. Long, truly-sequential tasks (AlfWorld) reward it with efficient, on-track
trajectories; short or shallow tasks (most of SearchQA, under a 4-turn cap) do not, and flat
outcome-only GRPO is competitive and cheaper. Where SearchQA \emph{is} deep (MuSiQue 3--4 hop), the
critic methods leave the faintest positive trace --- but the 4-turn cap, not the credit signal, is
the binding constraint (the actionable fix there is more turns, not better credit).

\textbf{Limitations.} (1) Mixed scale (AlfWorld 1.7B, SearchQA 4B) --- we compare \emph{patterns},
not levels. (2) AlfWorld is single-seed, with unequal training budgets (both PPO arms 150 steps vs 300 for
GRPO/GiGPO); the turn-vs-token contrast is budget-matched at 150, but the PPO-vs-GRPO/GiGPO
comparisons are not. (3) $H_{\text{gold}}$ is annotated chain length, ordinal, not ``expected
searches''; the 3-hop stratum is MuSiQue-only. (4) SearchQA metric is greedy val-protocol EM on the
full test set. (5) Fig 5's accuracy-vs-turns axis is partly mechanical (hard items take more turns).

\vspace{4pt}\noindent\footnotesize Assets, rerunnable analysis, and $H_{\text{gold}}$ table:
\texttt{docs/reports/2026-07-21\_4b\_horizon\_eval/} (see \texttt{README.md}).
\end{document}
